\documentclass[a4paper,11pt]{article}
\usepackage[utf8]{inputenc}
\usepackage[T1]{fontenc}
\usepackage[ngerman]{babel}
\usepackage[margin=2cm, top=2.2cm, bottom=2.2cm]{geometry}
\usepackage{enumitem}
\usepackage{booktabs}
\usepackage{tabularx}
\usepackage{xcolor}
\usepackage{fancyhdr}
\usepackage{titlesec}
\usepackage{listings}
\usepackage{mdframed}
\usepackage{parskip}

% ── Farben ──
\definecolor{beezubi}{RGB}{0, 90, 160}
\definecolor{grau}{RGB}{100, 100, 100}
\definecolor{codebg}{RGB}{245, 245, 245}
\definecolor{codeframe}{RGB}{180, 180, 180}
\definecolor{hinweis}{RGB}{255, 243, 205}
\definecolor{hinweisrand}{RGB}{200, 160, 0}
\definecolor{erfolg}{RGB}{220, 245, 220}
\definecolor{erfolgrand}{RGB}{0, 140, 0}

% ── Kopfzeile ──
\pagestyle{fancy}
\fancyhf{}
\fancyhead[L]{\small\color{beezubi}\textbf{BeeZubi Lernwelt GmbH}}
\fancyhead[R]{\small\color{grau}Kurzanleitung — \texttt{csv\_translator.py}}
\fancyfoot[C]{\small\thepage}
\renewcommand{\headrulewidth}{0.4pt}

% ── Überschriften ──
\titleformat{\section}{\large\bfseries\color{beezubi}}{}{0em}{}[\vspace{-0.2em}\rule{\linewidth}{0.4pt}]
\titleformat{\subsection}{\normalsize\bfseries\color{beezubi}}{}{0em}{}
\titlespacing*{\section}{0pt}{1.2ex}{0.6ex}
\titlespacing*{\subsection}{0pt}{0.8ex}{0.3ex}

% ── Code ──
\lstset{
  backgroundcolor=\color{codebg},
  frame=single, rulecolor=\color{codeframe},
  basicstyle=\small\ttfamily,
  breaklines=true, keepspaces=true,
  xleftmargin=6pt, xrightmargin=6pt,
  aboveskip=4pt, belowskip=4pt,
}

% ── Boxen ──
\newenvironment{hinweisbox}{%
  \begin{mdframed}[backgroundcolor=hinweis,linecolor=hinweisrand,linewidth=1pt,
    innerleftmargin=8pt,innerrightmargin=8pt,innertopmargin=5pt,innerbottommargin=5pt]
  \small}{\end{mdframed}}

\newenvironment{erfolgbox}{%
  \begin{mdframed}[backgroundcolor=erfolg,linecolor=erfolgrand,linewidth=1pt,
    innerleftmargin=8pt,innerrightmargin=8pt,innertopmargin=5pt,innerbottommargin=5pt]
  \small}{\end{mdframed}}

\setlength{\parindent}{0pt}

% ────────────────────────────────────────────────────────────────────
\begin{document}

\begin{center}
  {\LARGE\bfseries\color{beezubi} Kurzanleitung}\\[0.3em]
  {\large CSV-Übersetzungsprogramm mit Parallelisierung}\\[0.2em]
  {\small\color{grau}\texttt{csv\_translator.py} \quad|\quad Stand: März 2026}
\end{center}

\vspace{0.3em}

% ── Voraussetzungen ──────────────────────────────────────────────────
\section{Voraussetzungen}

\begin{enumerate}[itemsep=3pt, leftmargin=1.5em]
  \item \textbf{Python 3.10+} installiert \quad
        \texttt{python -{}-version}
  \item \textbf{Pakete installiert:}
\begin{lstlisting}
pip install pandas python-docx openai ftfy
\end{lstlisting}
  \item \textbf{API-Schlüssel} in der Datei \texttt{.env} im Projektordner:
\begin{lstlisting}
OPENAI_API_KEY=sk-proj-IhrSchluesselHier
\end{lstlisting}
  \item \textbf{Eingabedateien} im Projektordner vorhanden:
        CSV-Datei + Megaprompt-DOCX-Datei
\end{enumerate}

% ── Standardaufruf ───────────────────────────────────────────────────
\section{Standardaufruf (sequentiell, eine Sprache nach der anderen)}

\begin{lstlisting}
python csv_translator.py `
    --pdf MeineDatei.csv `
    --prompt MeinMegaprompt.docx `
    --encoding utf-8-sig `
    --sep ";" `
    --langage AR TR RU ES PL EN RO BG `
    --batch-size 200 `
    --pruefung `
    --column-order "FrageNr;lfdNr;BerufNr;Beruf;Abschlussprüfung Teil 1;Abschlussprüfung Teil 2;Lehrjahr;LFNr;LF;AbschnNr;Abschnitt;Nr;Frage;A;B;C;D;Richtig1;Richtig_Text1;Bild;Sprache"
\end{lstlisting}

% ── Parallelaufruf ───────────────────────────────────────────────────
\section{Parallelaufruf (2 Sprachen gleichzeitig — empfohlen)}

Fügen Sie einfach \texttt{-{}-parallel 2} am Ende hinzu:

\begin{lstlisting}
python csv_translator.py `
    --pdf MeineDatei.csv `
    --prompt MeinMegaprompt.docx `
    --encoding utf-8-sig `
    --sep ";" `
    --langage AR TR RU ES PL EN RO BG `
    --batch-size 200 `
    --pruefung `
    --column-order "FrageNr;lfdNr;BerufNr;Beruf;Abschlussprüfung Teil 1;Abschlussprüfung Teil 2;Lehrjahr;LFNr;LF;AbschnNr;Abschnitt;Nr;Frage;A;B;C;D;Richtig1;Richtig_Text1;Bild;Sprache" `
    --parallel 2
\end{lstlisting}

\begin{erfolgbox}
\textbf{Zeitersparnis:} Mit \texttt{-{}-parallel 2} laufen immer 2 Sprachen gleichzeitig.
Bei 8 Sprachen à 30 Minuten: statt 4h nur noch \textasciitilde2h.\\[3pt]
\textbf{Rate Limits:} Kein Risiko — \texttt{gpt-5-mini} erlaubt 10.000.000 Tokens/Minute.
Der parallele Betrieb nutzt davon weniger als 5\%.
\end{erfolgbox}

% ── Parameter ────────────────────────────────────────────────────────
\section{Wichtigste Parameter}

\begin{tabularx}{\textwidth}{@{}l l X@{}}
\toprule
\textbf{Parameter} & \textbf{Standard} & \textbf{Beschreibung} \\
\midrule
\texttt{-{}-pdf} & \textit{Pflicht} & Eingabe-CSV-Datei \\[3pt]
\texttt{-{}-prompt} & \textit{Pflicht} & Word-Datei mit Übersetzungsregeln \\[3pt]
\texttt{-{}-encoding} & \texttt{utf-8} & \texttt{utf-8-sig} für Excel-Dateien \\[3pt]
\texttt{-{}-sep} & \textit{auto} & CSV-Trennzeichen (\texttt{;} oder \texttt{,}) \\[3pt]
\texttt{-{}-langage} & alle 8 Sprachen & Zielsprachen, durch Leerzeichen getrennt \\[3pt]
\texttt{-{}-batch-size} & \texttt{200} & Zeilen pro API-Aufruf \\[3pt]
\texttt{-{}-pruefung} & \textit{inaktiv} & Benennt \texttt{Zwischenprüfung} → \texttt{Abschlussprüfung Teil 1} um \\[3pt]
\texttt{-{}-column-order} & \textit{keine} & Spalten der Ausgabe (nur diese werden ausgegeben) \\[3pt]
\texttt{-{}-parallel} & \texttt{1} & Anzahl gleichzeitiger Sprachen \\
\bottomrule
\end{tabularx}

% ── Sprachcodes ──────────────────────────────────────────────────────
\section{Unterstützte Sprachen}

\begin{tabular}{@{}llllllll@{}}
\textbf{AR} Arabisch &
\textbf{TR} Türkisch &
\textbf{RU} Russisch &
\textbf{ES} Spanisch &
\textbf{PL} Polnisch &
\textbf{EN} Englisch &
\textbf{RO} Rumänisch &
\textbf{BG} Bulgarisch \\
\end{tabular}

% ── Ergebnisse ───────────────────────────────────────────────────────
\section{Ergebnisse finden}

Alle Ausgabedateien befinden sich im Ordner \texttt{out\_translated\_backofff/}:

\begin{tabular}{@{}ll@{}}
\texttt{MeineDatei\_AR.csv} & Übersetzte Datei Arabisch \\
\texttt{MeineDatei\_TR.csv} & Übersetzte Datei Türkisch \\
\texttt{...} & (eine Datei pro Sprache) \\
\texttt{protokoll\_MeineDatei.csv} & Statistiken aller Läufe \\
\end{tabular}

% ── Unterbrechen ─────────────────────────────────────────────────────
\section{Unterbrechen und fortsetzen}

Das Programm kann jederzeit mit \textbf{Strg+C} gestoppt werden.
Beim Neustart wird automatisch an der letzten gespeicherten Position fortgesetzt —
für jede Sprache unabhängig.

\begin{hinweisbox}
\textbf{CSV in Excel geöffnet?} → Schließen Sie Excel zuerst, sonst erscheint ein
\texttt{PermissionError}. Das Programm wartet automatisch 5 Sekunden und versucht es erneut.
\end{hinweisbox}

\vspace{1em}
\begin{center}
\rule{0.5\linewidth}{0.4pt}\\[0.3em]
{\small\color{grau}BeeZubi Lernwelt GmbH — März 2026}
\end{center}

\end{document}
